\section{多自由度小振动与简正模（选读）}\label{ch10:sec:modes}
\subsection{惯性矩阵和弹性矩阵从哪里来}
一个单摆只需一个角度，但双摆需要两个角度。具有$n$个独立广义坐标的系统，在平衡构形附近令$q_i$表示第$i$个坐标的微小偏移，并组成列向量$\vect q$。在光滑势能、无耗散、约束不显含时间的条件下，平衡使一次项消失，二阶近似为
\begin{equation}\label{ch10:eq:quadratic}
 T_{\rm kin}=\frac12\sum_{i,j}M_{ij}\dot q_i\dot q_j
 =\frac12\dot{\vect q}^{\mathsf T}\mathbf M\dot{\vect q},\qquad
 U-U_0=\frac12\sum_{i,j}K_{ij}q_iq_j
 =\frac12\vect q^{\mathsf T}\mathbf K\vect q.
\end{equation}
上标$\mathsf T$表示转置；$\mathbf M$是惯性矩阵，$\mathbf K$是弹性矩阵。其元素分别为动能对速度的二阶偏导和势能在平衡点的二阶偏导。混合偏导可交换，故两矩阵对称。即使最初把交叉项写成不对称系数，也可替换为对称部分，因为$\vect q^{\mathsf T}(\mathbf B-\mathbf B^{\mathsf T})\vect q=0$。

例如$U=aq_1^2/2+cq_1q_2+dq_2^2/2$的矩阵是$\bigl(\begin{smallmatrix}a&c\\c&d\end{smallmatrix}\bigr)$，不是把两个非对角元都写成$2c$。求$\partial U/\partial q_i$得到$\sum_jK_{ij}q_j$，所以拉格朗日方程给出
\begin{equation}\label{ch10:eq:matrix-motion}
 \mathbf M\ddot{\vect q}+\mathbf K\vect q=0.
\end{equation}

对于真正独立、非奇异的物理坐标，非零广义速度应具有正动能，即$\mathbf M$正定。若出现零惯性方向，要检查是否引入了冗余坐标、忽略了必要惯量，或在该平衡点使用了奇异坐标；不能直接求逆继续套用。势能矩阵则有三种重要情况：正定保证严格局部稳定；出现负方向意味着不稳定；有零方向时需看对称性或更高阶势能。零频整体转动可以是中性运动，$U\propto q^4$也可能稳定，而$U\propto-q^4$不稳定。

\subsection{两个模型：耦合摆与双摆}
先看图\ref{ch10:fig:coupled-pendula}：两根摆长均为$l$，摆球均为$m$，固定悬点水平相距$d$，摆球间轻弹簧原长恰为$d$、劲度系数为$k_c$。两角都从各自竖直向下方向向右量为正，不以图中左右位置改变角度正向。
\chTenCoupledPendulaFigure
两球坐标差的水平分量为$d+l\sin\theta_2-l\sin\theta_1$，竖直分量为$l\cos\theta_1-l\cos\theta_2$，故弹簧长度
\[
 d'=\sqrt{(d+l\sin\theta_2-l\sin\theta_1)^2+
 l^2(\cos\theta_1-\cos\theta_2)^2}.
\]
小角下$d'-d=l(\theta_2-\theta_1)+O(\theta^2)$，于是
\begin{align*}
 T_{\rm kin}&=\frac12ml^2(\dot\theta_1^2+\dot\theta_2^2),\\
 U&=mgl(2-\cos\theta_1-\cos\theta_2)+\frac12k_c(d'-d)^2\\
 &\simeq\frac12mgl(\theta_1^2+\theta_2^2)
 +\frac12k_cl^2(\theta_2-\theta_1)^2.
\end{align*}
记$\epsilon=k_cl/(mg)$，得
\[
 \mathbf M=ml^2\begin{pmatrix}1&0\\0&1\end{pmatrix},\qquad
 \mathbf K=mgl\begin{pmatrix}1+\epsilon&-\epsilon\\-\epsilon&1+\epsilon\end{pmatrix}.
\]
这里耦合出现在势能中。交换两摆不改变方程，故先试同相$\theta_1=\theta_2$与反相$\theta_1=-\theta_2$：
\[
 \ddot\theta_1=-\frac gl[(1+\epsilon)\theta_1-\epsilon\theta_2],
 \quad\omega_S^2=\frac gl,
 \qquad\omega_A^2=\frac gl(1+2\epsilon).
\]
同相时弹簧一阶伸长为零；反相时两端位移反向，伸长加倍。取$Q_S=(\theta_1+\theta_2)/2$、$Q_A=(\theta_1-\theta_2)/2$，直接相加、相减方程也得到两条独立简谐方程。任意初始状态是这两种模式的叠加，不必始终具有左右对称形状。

同一数学结构也适用于两物块弹簧模型，见图\ref{ch10:fig:coupled-masses}。两端弹簧均为$k$，中间弹簧为$k_c$，$x_1,x_2$从各自平衡位置向右量：
\[
 m\ddot x_1=-(k+k_c)x_1+k_cx_2,\qquad
 m\ddot x_2=k_cx_1-(k+k_c)x_2.
\]
因此同相、反相频率分别为$\sqrt{k/m}$和$\sqrt{(k+2k_c)/m}$。中间弹簧的伸长量为$x_2-x_1$，不是$x_1+x_2$。
\chTenCoupledMassFigure

再看图\ref{ch10:fig:double-pendulum}的等质量、等杆长双摆。两角均相对竖直方向，而不是第二杆相对第一杆的夹角。上球坐标$(l\sin\theta_1,-l\cos\theta_1)$，下球再加$(l\sin\theta_2,-l\cos\theta_2)$，故下球速率平方为
\[
 v_2^2=l^2\dot\theta_1^2+l^2\dot\theta_2^2
 +2l^2\dot\theta_1\dot\theta_2\cos(\theta_1-\theta_2).
\]
\chTenDoublePendulumFigure
将两球动能相加，重力势能相加，小角近似给
\begin{align*}
 T_{\rm kin}&\simeq\frac12ml^2(2\dot\theta_1^2+2\dot\theta_1\dot\theta_2+\dot\theta_2^2),\\
 U&=mgl[2(1-\cos\theta_1)+(1-\cos\theta_2)]
 \simeq\frac12mgl(2\theta_1^2+\theta_2^2),\\
 \mathbf M&=ml^2\begin{pmatrix}2&1\\1&1\end{pmatrix},\qquad
 \mathbf K=mgl\begin{pmatrix}2&0\\0&1\end{pmatrix}.
\end{align*}
这里势能没有交叉项，耦合却在动能中。仅把一种能量配方并不够。例如令$u=\theta_1+\theta_2/2,w=\theta_2$，则
\[
 T_{\rm kin}=ml^2\dot u^2+\frac14ml^2\dot w^2,
 \quad U=mgl\left(u^2-uw+\frac34w^2\right).
\]
动能虽已化为平方和，势能反而出现交叉项；真正目标是用\emph{同一个}坐标变换同时对角化两者。

\subsection{同时对角化、正交性和简正坐标}
一般二次型可以逐变量配方。例如$f_{11}\ne0$时，把$q_1$换成$q_1'=q_1+\sum_{j>1}f_{1j}q_j/f_{11}$，便可将含$q_1$的项写成$f_{11}(q_1')^2/2$和其余变量的二次型；继续处理余项即可消去交叉项。零主元时需要先换坐标。这个过程叫二次型的对角化，但它未保证同时处理两个二次型。

令旧坐标$\vect q=\mathbf B\vect Q$。代入式\eqref{ch10:eq:quadratic}，得到新矩阵
\[
 \mathbf M'=\mathbf B^{\mathsf T}\mathbf M\mathbf B,
 \qquad\mathbf K'=\mathbf B^{\mathsf T}\mathbf K\mathbf B.
\]
这称合同变换，与$\mathbf B^{-1}\mathbf K\mathbf B$那种相似变换不能混淆。如何选$\mathbf B$？先设所有坐标以同频率运动，$\vect q=\vect a\cos(\omega t+\phi)$，代入动力学方程：
\begin{equation}\label{ch10:eq:eigen}
 (\mathbf K-\omega^2\mathbf M)\vect a=0,
 \qquad\det(\mathbf K-\omega^2\mathbf M)=0.
\end{equation}
前式是广义本征值问题，后式是存在非零振幅向量的条件。但“行列式有$n$个根”还不足以保证$n$个独立本征向量，下面补上保证完备性的论证。

$\mathbf M$实对称正定，故有实对称正定平方根$\mathbf M^{1/2}$和逆$\mathbf M^{-1/2}$。令$\vect u=\mathbf M^{1/2}\vect a$，广义本征方程化为
\[
 \underbrace{\mathbf M^{-1/2}\mathbf K\mathbf M^{-1/2}}_{\mathbf D}
 \vect u=\lambda\vect u,\qquad\lambda=\omega^2.
\]
$\mathbf D$为实对称矩阵，所以存在$n$个实特征值及一组完备正交特征向量；重根子空间内也能选择正交基。令这些单位向量为$\vect u_i$，取$\vect a_i=\mathbf M^{-1/2}\vect u_i$，便有
\[
 \vect a_i^{\mathsf T}\mathbf M\vect a_j=\delta_{ij},\qquad
 \vect a_i^{\mathsf T}\mathbf K\vect a_j=\lambda_j\delta_{ij}.
\]
这既证明存在性，也给出构造。不同频率的质量正交性还可直接验证：由$\mathbf K\vect a_j=\lambda_j\mathbf M\vect a_j$及矩阵对称性，
\[
 \lambda_j\vect a_i^{\mathsf T}\mathbf M\vect a_j
 =\vect a_i^{\mathsf T}\mathbf K\vect a_j
 =\lambda_i\vect a_i^{\mathsf T}\mathbf M\vect a_j,
\]
所以$\lambda_i\ne\lambda_j$时该内积为零。若不把向量归一化，可记$M_i=\vect a_i^{\mathsf T}\mathbf M\vect a_i>0$，$K_i=\lambda_iM_i$；取$\mathbf B=(\vect a_1,\ldots,\vect a_n)$后，
\[
 T_{\rm kin}=\frac12\sum_i M_i\dot Q_i^2,
 \qquad U-U_0=\frac12\sum_i K_iQ_i^2,
 \qquad M_i\ddot Q_i+K_iQ_i=0.
\]
$Q_i$称简正坐标，$\omega_i=\sqrt{K_i/M_i}$称简正频率；只有一个$Q_i$被激发的运动称简正模（normal mode）。$\lambda_i>0$时$Q_i=C_i\cos(\omega_it+\phi_i)$；$\lambda_i=0$时$Q_i=C_i+D_it$；$\lambda_i<0$时出现增长/衰减指数。因此“$\mathbf M$正定、$\mathbf K$对称”只保证能解耦，并不保证每个坐标都振动。

初始条件也可明确投影：$Q_i(0)=\vect a_i^{\mathsf T}\mathbf M\vect q(0)/M_i$，$\dot Q_i(0)=\vect a_i^{\mathsf T}\mathbf M\dot{\vect q}(0)/M_i$。分别解每个标量方程，再用$\vect q=\sum_i\vect a_iQ_i$还原运动。原坐标通常含多个频率，并不各自简谐。

\subsection{完整求出双摆的两个模态}
在式\eqref{ch10:eq:eigen}中提出公共因子，令$\alpha=g/l$、$\lambda=\omega^2$，得到
\[
 \det\begin{pmatrix}2\alpha-2\lambda&-\lambda\\-\lambda&\alpha-\lambda\end{pmatrix}
 =2(\alpha-\lambda)^2-\lambda^2=0.
\]
展开为$\lambda^2-4\alpha\lambda+2\alpha^2=0$，故
\[
 \lambda_-=(2-\sqrt2)g/l,\qquad
 \lambda_+=(2+\sqrt2)g/l.
\]
第二行给$a_2/a_1=\lambda/(\alpha-\lambda)$，分别为$\sqrt2$和$-\sqrt2$。可选$\vect a_-=(1,\sqrt2)^{\mathsf T}$、$\vect a_+=(1,-\sqrt2)^{\mathsf T}$。所以同向摆动的低频模中，下杆角振幅是上杆的$\sqrt2$倍；高频模两杆反向，绝对比仍为$\sqrt2$，不是简单的等幅反相。

令
\[
 \theta_1=Q_-+Q_+,\quad\theta_2=\sqrt2(Q_--Q_+),
 \qquad Q_\pm=\frac{\theta_1}{2}\mp\frac{\theta_2}{2\sqrt2}.
\]
这里最后一式中$Q_+$取减号、$Q_-$取加号。代入可验算
\begin{align*}
 T_{\rm kin}&=\frac12[2(2+\sqrt2)ml^2]\dot Q_-^2
 +\frac12[2(2-\sqrt2)ml^2]\dot Q_+^2,\\
 U&=\frac12(4mgl)Q_-^2+\frac12(4mgl)Q_+^2.
\end{align*}
分别用弹性系数除惯性系数，恰好得到上面的$\lambda_\pm$，完成独立校验。图\ref{ch10:fig:double-modes}画的是各模的一个瞬时构形，杆长始终为$l$，箭头表示角运动的相对方向而不是独立平移。
\chTenDoubleModesFigure

\subsection{六原子环：从对称操作到六个模式}
图\ref{ch10:fig:benzene-model}左为苯分子的立体模型，右为简化的六原子环。真实苯含六个碳和六个氢；在右图的简化中，每个碳氢单元只当作质量$m$的质点，被限制在固定圆环上作切向小位移$x_j$，相邻单元用有效劲度$k$耦合。它只保留六个坐标，不能代替全部分子振动。
\chTenBenzeneModelFigure
给六点按逆时针编号$j=0,\ldots,5$，所有$x_j$的正向均为逆时针切向，循环约定$x_{j+6}=x_j$。有效能量和方程为
\[
 T_{\rm kin}=\frac12m\sum_{j=0}^5\dot x_j^2,\qquad
 U=\frac12k\sum_{j=0}^5(x_{j+1}-x_j)^2,
 \qquad m\ddot x_j=k(x_{j+1}+x_{j-1}-2x_j).
\]
这里$k$定义的是相邻切向坐标差的弹性系数，不把画在弦上的普通伸缩弹簧刚度不加换算地当成$k$。

对称操作的意义是使能量函数和动力学规律不变，而非要求任意时刻的实际位形都对称。图\ref{ch10:fig:ring-symmetry}依次画出绕环轴转$60^\circ$的$C_6$、绕对边中点连线转$180^\circ$的$C_2$和绕对顶点连线转$180^\circ$的$C'_2$。后两种在环面投影中表现为反射，切向正向也随之反转。以图中轴为例，其作用可写为
\[
 (C_6x)_j=x_{j-1},\qquad (C_2x)_j=-x_{1-j},\qquad
 (C'_2x)_j=-x_{-j},
\]
下标按模6理解。每个操作只是对$\sum\dot x_j^2$、$\sum(x_{j+1}-x_j)^2$重新编号，所以能量不变；单独交换任意两点一般会改变相邻关系，因而不一定是对称操作。
\chTenRingSymmetryFigure
如果$x_j(t)$是一个解，对它作对称操作仍得到解且不改变频率。初始条件却可以不对称，因此一般解不必保持某种对称形状；应区分“方程的对称性”和“所选状态的对称性”，不能把一切不对称初值都等同于严格意义的自发对称性破缺。

设$\mathbf C=C_6$，有$\mathbf C^6=\mathbf I$，其本征值只能是六次单位根。对任意向量$\vect a$，令$w=\mathrm e^{2\pi\ii/6}$，构造
\[
 \vect a^{(n)}=\sum_{s=0}^5w^{-ns}\mathbf C^s\vect a,
 \qquad\mathbf C\vect a^{(n)}=w^n\vect a^{(n)}.
\]
等号由把求和指标平移一位并用$w^{6n}=1$得到；某个投影可以为零，不能把零向量当模式。由于动力学矩阵与$\mathbf C$可交换，同一频率的子空间在$\mathbf C$作用下保持不变，可在其中选择这些投影。反射操作的平方为恒等，其本征值则只能为$\pm1$。

更直接地，设$x_j=A\exp[\ii(2\pi nj/6-\omega t)]$，代入每个质点的方程，约去公共因子得
\[
 -m\omega_n^2=k(\mathrm e^{2\pi\ii n/6}+\mathrm e^{-2\pi\ii n/6}-2),
 \quad\omega_n^2=\frac{4k}{m}\sin^2\frac{n\pi}{6}.
\]
所以频率平方为$0,k/m,3k/m,4k/m$，简并度依次为1、2、2、1。$n$与$6-n$是一对共轭行波，实部和虚部给两个独立实驻波模式。零频模是整体沿环匀速转动，不是周期无穷大的普通正弦振动。

也可不用先猜行波，而从两种二重轴的正负特征构造实模式。为方便分类，在对顶点轴这一类中选穿过$j=2,5$的轴，记其操作为$D_2$，于是$(D_2x)_j=-x_{4-j}$。它与上面的$C_2$轴垂直，两操作可交换。联立$C_2\vect x=\sigma\vect x$和$D_2\vect x=\tau\vect x$，其中$\sigma,\tau=\pm1$，可以选出
\begin{center}\small
\begin{tabular}{ccl}
\toprule $\sigma$&$\tau$&一个代表向量\\\midrule
$-1$&$-1$&$(1,1,1,1,1,1)$\\
$+1$&$-1$&$(1,-1,1,-1,1,-1)$\\
$-1$&$+1$&$(1,1,0,-1,-1,0)$\\
$+1$&$+1$&$(1,-1,0,1,-1,0)$\\\bottomrule
\end{tabular}\end{center}
将第一分量代回$m\ddot x_0=k(x_1+x_5-2x_0)$：均匀向量给右端为零；交替向量给$-4kx_0$；后两个向量分别给$-kx_0$与$-3kx_0$。其余分量也须代入验证，例如后两种的静止点$j=2,5$，两邻点位移相反，合回复力确为零，因而可以一直静止。

前两种向量各自也是$C_6$的本征向量，分别给本征值$+1$和$-1$。后两种经$C_6$作用却变成不成比例的新向量，仍有同一频率，故必须补入各自的简并伙伴。单凭二重轴的两个正负号，不能把每个代表态都当成独立的单重频率。

用$+$、$-$、0表示相同绝对值的正、负和零位移，下列六个实向量张成全部坐标空间：
\begin{center}\small
\begin{tabular}{cll}
\toprule
$m\omega^2/k$&代表向量&独立模式数\\\midrule
0&$(1,1,1,1,1,1)$&1\\
1&$(1,1,0,-1,-1,0)$，$(0,1,1,0,-1,-1)$&2\\
3&$(1,-1,0,1,-1,0)$，$(0,1,-1,0,1,-1)$&2\\
4&$(1,-1,1,-1,1,-1)$&1\\\bottomrule
\end{tabular}\end{center}
这组简并向量可再正交化，但不是必须先正交才能验证模式。两种二重模的代表向量分别满足
\[
 \begin{aligned}
 \vect a&=(1,1,0,-1,-1,0),&\mathbf C^2\vect a&=\mathbf C\vect a-\vect a,\\
 \vect b&=(1,-1,0,1,-1,0),&\mathbf C^2\vect b&=-\mathbf C\vect b-\vect b.
 \end{aligned}
\]
两种关系符号不同，分别对应单位根实部$1/2$与$-1/2$。图\ref{ch10:fig:ring-modes}画出两组代表态及旋转后的伙伴，箭头严格沿所在点的圆周切线；不动点明确保留。
\chTenRingModesFigure
交替向量经$C_6$后变号，故其$C_6$本征值为$-1$；均匀向量的本征值为$+1$。对于二重简并，某一个实向量通常会被旋转为同频子空间中的另一向量，而不是自身乘一个实数。对称操作的乘法结构及其在线性空间中的作用，分别属于群论和群表示论；上述六环是它们在分子与固体振动中的一个具体入口。

\subsection{真实分子的振型与光谱}
一个非线性的12原子分子有$3\times12=36$个位置自由度。除去3个整体平动和3个整体转动，剩30个内部振动自由度。苯的对称性使其中10对两两简并，故可分成20组基频。图\ref{ch10:fig:benzene-modes-a}至图\ref{ch10:fig:benzene-modes-d}逐组列出相应振型；每个双重组有两幅独立图，不能把30个振动自由度误数成20个。

读图时，深色大圆为碳，浅色小圆为氢；平面内箭头表示一个选定相位的主要位移方向，$\odot$表示出纸面，$\otimes$表示入纸面。半个周期后所有位移同时反向。这里给出振型的定性示意，箭头不按本征矢量的分量比例绘制，不能用箭头长度计算质心位移或动能；精确振型还须由力常数矩阵求出。波数$\widetilde\nu$的单位为$\mathrm{cm^{-1}}$，与频率的关系为$\nu=c\widetilde\nu$、$\omega=2\pi\nu$；模型参数、同位素和测量条件会影响数值，不能把图上近似基频当成无误差常数。A、B表示单重，E表示二重；g、u表示相应反演对称性，进一步的光谱选择定则需结合电偶极矩或极化率变化。
图\ref{ch10:fig:benzene-modes-a}中的点、叉可用来区分相邻原子的同向和反向离面运动；图\ref{ch10:fig:benzene-modes-b}中，$1014\,\mathrm{cm^{-1}}$一组的各原子同时向环心收拢，而$994\,\mathrm{cm^{-1}}$一组的相邻单元交替伸缩。图\ref{ch10:fig:benzene-modes-c}中$1332\,\mathrm{cm^{-1}}$一组的碳、氢切向运动相反，不能看成整个分子刚性转动。图\ref{ch10:fig:benzene-modes-d}中较高频的位移以碳氢键方向上的伸缩为主；双重组的两幅图应分别看作两个独立振动坐标。
% Source figure 4.11, physical page 93: twenty groups, thirty panels.
% Arrows are qualitative normal-pattern directions, not computed amplitudes.
\newcommand{\chTenMolecule}[4]{%
\begin{tikzpicture}[scale=.55,>=Latex]
\foreach \j in {0,...,5}{\draw[TextGray]({60*\j}:1)--({60*(\j+1)}:1);\draw[TextGray]({60*\j}:1)--({60*\j}:1.7);}
\foreach \j in {0,...,5}{\draw[DeepBlue,thick]({60*\j}:1)circle(.10);\draw[TextGray]({60*\j}:1.7)circle(.075);}
\ifnum#1=0
 \foreach \v[count=\j from 0] in {#2}{\ifnum\v=0\else
 \node[font=\tiny,inner sep=0,fill=none]at({60*\j}:1){\ifnum\v>0$\cdot$\else$\times$\fi};\fi}
 \foreach \v[count=\j from 0] in {#3}{\ifnum\v=0\else
 \node[font=\tiny,inner sep=0,fill=none]at({60*\j}:1.7){\ifnum\v>0$\cdot$\else$\times$\fi};\fi}
\else
 \foreach \ang/\mag[count=\j from 0] in {#2}{\ifdim\mag pt=0pt\else\draw[DeepBlue,-{Latex[length=.8mm]},line width=.45pt]({60*\j}:1)++(\ang:.12)--++(\ang:{.45*max(.6,\mag)});\fi}
 \foreach \ang/\mag[count=\j from 0] in {#3}{\ifdim\mag pt=0pt\else\draw[DeepBlue,-{Latex[length=.8mm]},line width=.45pt]({60*\j}:1.7)++(\ang:.09)--++(\ang:{.42*max(.6,\mag)});\fi}
\fi
\path[use as bounding box](-2.5,-2.4)rectangle(2.5,2.4);
\end{tikzpicture}%
}
\newcommand{\chTenModeGroup}[2]{\begin{minipage}[t]{.47\linewidth}\centering #2\par\smallskip{\small #1}\end{minipage}}
\newcommand{\chTenZeroVectors}{0/0,0/0,0/0,0/0,0/0,0/0}
\begin{chTenFigure}
\chTenModeGroup{$395\,\mathrm{cm^{-1}}\ (E_{2u})$}{
\chTenMolecule{0}{-1,1,1,-1,1,1}{-1,1,1,-1,1,1}{}
\chTenMolecule{0}{1,1,-1,1,1,-1}{1,1,-1,1,1,-1}{}}
\hfill
\chTenModeGroup{$599\,\mathrm{cm^{-1}}\ (E_{2g})$}{
\chTenMolecule{1}{0/1,240/.6,300/.6,180/1,60/.6,120/.6}{0/1,240/.6,300/.6,180/1,60/.6,120/.6}{}
\chTenMolecule{1}{0/0,60/1,300/1,0/0,240/1,120/1}{0/0,90/1,270/1,0/0,270/1,90/1}{}}
\par\medskip
\chTenModeGroup{$677\,\mathrm{cm^{-1}}\ (A_{2u})$}{
\chTenMolecule{0}{1,1,1,1,1,1}{-1,-1,-1,-1,-1,-1}{}}
\hfill
\chTenModeGroup{$700\,\mathrm{cm^{-1}}\ (B_{2g})$}{
\chTenMolecule{0}{1,-1,1,-1,1,-1}{1,-1,1,-1,1,-1}{}}
\par\medskip
\chTenModeGroup{$838\,\mathrm{cm^{-1}}\ (E_{1g})$}{
\chTenMolecule{0}{-1,-1,-1,1,1,1}{1,1,1,-1,-1,-1}{}
\chTenMolecule{0}{1,1,-1,-1,-1,1}{-1,-1,1,1,1,-1}{}}
\figcap{苯分子振型（一）：五组、八个独立模式。离面模式用圆内点、叉标示；未标箭头的方向不表示该原子永远静止。}\label{ch10:fig:benzene-modes-a}
\end{chTenFigure}

\begin{chTenFigure}
\chTenModeGroup{$954\,\mathrm{cm^{-1}}\ (E_{2u})$}{
\chTenMolecule{0}{-1,1,-1,-1,1,-1}{1,-1,1,1,-1,1}{}
\chTenMolecule{0}{-1,1,1,-1,1,1}{1,-1,-1,1,-1,-1}{}}
\hfill
\chTenModeGroup{$981\,\mathrm{cm^{-1}}\ (B_{2g})$}{
\chTenMolecule{0}{1,-1,1,-1,1,-1}{-1,1,-1,1,-1,1}{}}
\par\medskip
\chTenModeGroup{$994\,\mathrm{cm^{-1}}\ (B_{2u})$}{
\chTenMolecule{1}{0/1,240/1,120/1,0/1,240/1,120/1}{0/1,240/1,120/1,0/1,240/1,120/1}{}}
\hfill
\chTenModeGroup{$1014\,\mathrm{cm^{-1}}\ (A_{1g})$}{
\chTenMolecule{1}{180/1,240/1,300/1,0/1,60/1,120/1}{180/.6,240/.6,300/.6,0/.6,60/.6,120/.6}{}}
\par\medskip
\chTenModeGroup{$1048\,\mathrm{cm^{-1}}\ (E_{1u})$}{
\chTenMolecule{1}{0/.7,0/0,0/0,0/.7,0/0,0/0}{180/1,315/1,45/1,180/1,315/1,45/1}{}
\chTenMolecule{1}{0/0,240/.6,300/.6,0/0,240/.6,300/.6}{90/1,180/1,0/1,90/1,180/1,0/1}{}}
\hfill
\chTenModeGroup{$1137\,\mathrm{cm^{-1}}\ (B_{1u})$}{
\chTenMolecule{1}{0/0,0/0,0/0,0/0,0/0,0/0}{90/1,330/1,210/1,90/1,330/1,210/1}{}}
\figcap{苯分子振型（二）：六组、八个独立模式。不同基频的同名对称类型仍是不同模。}\label{ch10:fig:benzene-modes-b}
\end{chTenFigure}

\begin{chTenFigure}
\chTenModeGroup{$1162\,\mathrm{cm^{-1}}\ (E_{2g})$}{
\chTenMolecule{1}{0/0,0/0,0/0,0/0,0/0,0/0}{0/0,330/1,210/1,0/0,150/1,30/1}{}
\chTenMolecule{1}{0/0,0/0,0/0,0/0,0/0,0/0}{270/1,150/1,150/1,90/1,330/1,330/1}{}}
\hfill
\chTenModeGroup{$1332\,\mathrm{cm^{-1}}\ (A_{2g})$}{
\chTenMolecule{1}{270/.12,330/.12,30/.12,90/.12,150/.12,210/.12}{90/1,150/1,210/1,270/1,330/1,30/1}{}}
\par\medskip
\chTenModeGroup{$1376\,\mathrm{cm^{-1}}\ (B_{1u})$}{
\chTenMolecule{1}{90/1,330/1,210/1,90/1,330/1,210/1}{0/0,0/0,0/0,0/0,0/0,0/0}{}}
\hfill
\chTenModeGroup{$1477\,\mathrm{cm^{-1}}\ (E_{1u})$}{
\chTenMolecule{1}{180/.4,180/1,0/1,180/.4,180/1,0/1}{180/1,315/1,45/1,180/1,315/1,45/1}{}
\chTenMolecule{1}{270/1,0/0,0/0,270/1,0/0,0/0}{90/1,0/0,0/0,90/1,0/0,0/0}{}}
\par\medskip
\chTenModeGroup{$1612\,\mathrm{cm^{-1}}\ (E_{2g})$}{
\chTenMolecule{1}{0/1,240/1,300/1,180/1,60/1,120/1}{270/1,0/0,180/1,90/1,0/0,0/1}{}
\chTenMolecule{1}{0/0,180/1,0/1,0/0,0/1,180/1}{0/0,150/1,30/1,0/0,330/1,210/1}{}}
\figcap{苯分子振型（三）：五组、八个独立模式。此组以平面内弯曲和环骨架形变为主。}\label{ch10:fig:benzene-modes-c}
\end{chTenFigure}

\begin{chTenFigure}
\chTenModeGroup{$3080\,\mathrm{cm^{-1}}\ (B_{2u})$}{
\chTenMolecule{1}{0/0,0/0,0/0,0/0,0/0,0/0}{0/1,240/1,120/1,0/1,240/1,120/1}{}}
\hfill
\chTenModeGroup{$3092\,\mathrm{cm^{-1}}\ (E_{2g})$}{
\chTenMolecule{1}{0/0,0/0,0/0,0/0,0/0,0/0}{0/0,240/1,120/1,0/0,60/1,300/1}{}
\chTenMolecule{1}{0/0,0/0,0/0,0/0,0/0,0/0}{0/1,240/.6,300/.6,180/1,60/.6,120/.6}{}}
\par\medskip
\chTenModeGroup{$3110\,\mathrm{cm^{-1}}\ (E_{1u})$}{
\chTenMolecule{1}{0/0,0/0,0/0,0/0,0/0,0/0}{0/0,240/1,300/1,0/0,240/1,300/1}{}
\chTenMolecule{1}{0/0,0/0,0/0,0/0,0/0,0/0}{0/1,0/0,0/0,0/1,0/0,0/0}{}}
\hfill
\chTenModeGroup{$3121\,\mathrm{cm^{-1}}\ (A_{1g})$}{
\chTenMolecule{1}{0/0,0/0,0/0,0/0,0/0,0/0}{0/1,60/1,120/1,180/1,240/1,300/1}{}}
\figcap{苯分子振型（四）：四组、六个独立模式，以碳氢键伸缩为主。四张图合计20组、30个模式；未画出的很小位移不可据此当成精确零。}\label{ch10:fig:benzene-modes-d}
\end{chTenFigure}

例如对称伸缩时，若分子的电偶极矩不随该坐标改变，即使外光频率相等，也未必产生电偶极红外吸收；有些模主要通过拉曼散射观察。频率相合是共振的条件之一，不是“每个模都必有同样强的红外谱线”。

经验势能或微观计算给出平衡点的力常数矩阵，经质量加权对角化便得到$\omega_i$和振型。量子描述中，每个近似独立模的能级为
\[
 E_{n_i}=\left(n_i+\frac12\right)\hbar\omega_i,\qquad n_i=0,1,2,\ldots.
\]
相邻能级相差$\hbar\omega_i$，故允许的单量子跃迁满足$h\nu=\hbar\omega_i$，即$\nu=\omega_i/(2\pi)$。这把经典简正频率与分子光谱联系起来；非谐性还会改变能级间隔并引入组合频率等现象。
